\section{引言}
\label{sec:introduction-zh}

异构图被广泛用于建模包含多种实体类型和多种关系类型的真实系统，例如学术网络、推荐平台、生物医学网络和社交媒体系统。与同构图相比，异构图能够通过关系类型和元路径提供更丰富的语义信息。这种语义结构使 HAN、HGT、MAGNN 和 HeCo 等异构图神经网络（HGNN）能够从多个语义邻域中聚合信息，从而学习更具表达能力的节点表示。这些模型已经在节点分类和表示学习任务上取得了较强性能。

尽管异构图神经网络有效，它们仍然容易受到对抗性结构扰动的影响。少量恶意插边或删边不仅会改变局部邻域，还会改变由元路径诱导出的高阶语义邻域。这个问题在异构图中尤其突出，因为某一种关系类型上的扰动可能在元路径组合之后被放大，并影响大量目标类型节点。因此，语义聚合过程可能传播误导性信息，注意力机制也可能给被污染的语义视图分配较大权重。

已有鲁棒图学习方法通常通过边剪枝、概率建模、图结构学习或基于相似性的消息传递来缓解结构攻击。近期的鲁棒异构图方法进一步考虑了类型关系和语义结构，例如估计边置信度，或者抑制可疑的元路径诱导连接。然而，大多数方法仍然从观测到的异构拓扑中构造候选邻域。当拓扑已经被严重污染时，仅依赖拓扑的防御方法恢复可靠表示的能力会受到限制，因为它们的去噪过程仍然受制于被攻击后的结构证据。

为了解决这一局限，本文提出 DVCL，即一种面向鲁棒异构图神经网络的双视图对比学习框架。DVCL 从元路径诱导的语义图中构造净化拓扑视图，并从目标节点特征相似性中构造特征诱导视图。拓扑视图通过元路径建模、语义注意力和两阶段边过滤保留异构语义表达能力；特征诱导视图则在结构攻击下提供互补的、相对独立于拓扑的参考。DVCL 进一步使用对称跨视图对比目标对齐两个视图，鼓励每个节点在语义拓扑和特征相似性之间保持一致表示，同时与其他节点保持可区分性。

本文主要贡献如下：
\begin{itemize}
    \item 提出 DVCL，一种双视图鲁棒异构图学习框架，联合利用净化语义拓扑和特征诱导相似性，在结构攻击下执行节点分类。
    \item 设计拓扑净化策略，将元路径级边置信度、语义注意力和全局语义过滤结合起来，并与特征诱导 KNN 视图配合使用，从而降低模型对被攻击图结构的依赖。
    \item 引入对称跨视图对比目标，并为干净图、全局攻击和目标攻击设置定义可复现评估协议，固定数据划分、攻击产物、目标节点并报告五随机种子结果。
\end{itemize}
